% Slot 2 (Thu 26 Nov 11:00--12:00 + Fri 27 Nov 09:30--10:15 Rome)

\section{Slot 2}

\begin{frame}{Slot 2 \textperiodcentered\ Thu 11:00--12:00 + Fri 09:30--10:15}
  \begin{center}
    \begin{tabular}{@{}rl@{}}
      \toprule
      \textbf{Clock} & \textbf{Block} \\
      \midrule
      Thu 11:00--11:25 & Skills: why \\
      11:25--11:55 & Create a Skill \\
      11:55--12:00 & Thursday wrap, overnight task \\
      \addlinespace[0.3em]
      Fri 09:30--09:35 & Recap \\
      09:35--09:55 & Sub-agent vs Skill \\
      09:55--10:15 & Rule of thumb: Skill or sub-agent? \\
      \bottomrule
    \end{tabular}
  \end{center}
\end{frame}

% ----------------------------------------------------------------
\begin{frame}{Skills: why}
  \timebox{Thu 11:00--11:25}
  \inplainwords{A Skill is a saved set of instructions for a task you do often.
    The agent follows it the same way every time, and you can share it with coauthors.}
  \likethis{a recipe card, in place of explaining the dish from memory each time.}
\end{frame}

\begin{frame}{Card--Krueger: \texttt{replication-checker}}
  Does our folder reproduce its numbers from a fresh start?
  \vspace{0.4em}
  \begin{itemize}
    \item File: \artifact{.claude/skills/replication-checker/SKILL.md}
    \item Inside: a name, when to use it, numbered steps.
    \item It runs the check command and confirms the data file is unchanged.
    \item It reports \textcolor{codegreen}{\textbf{GREEN}}, \textbf{YELLOW} or \textcolor{codered}{\textbf{RED}}, with the output as evidence.
    \item It never certifies causality.
  \end{itemize}
\end{frame}

\begin{frame}{Your turn}
  \youtry{Thu 11:00--11:25}
  \vspace{0.3em}
  \begin{enumerate}
    \item Read the Skill file before you trust it: what may it run?
    \item Run it on \artifact{examples/card-krueger-toy/}.
    \item Write down its colour and the one line of evidence behind it.
  \end{enumerate}
  \vspace{0.5em}
  {\small Guide: \artifact{demos/slot2\_sdd\_skill\_pr.md} \textperiodcentered\ \artifact{exercises/06\_slot2\_sdd\_skill\_pr.md}}
\end{frame}

% ----------------------------------------------------------------
\begin{frame}{Create a Skill}
  \timebox{Thu 11:25--11:55}
  \inplainwords{Make a folder and write a short file in it with a name, a line saying
    when to use it, and numbered steps. Test it once and fix whatever was unclear.}
  \likethis{writing down your referee-report checklist so a student can follow it.}
\end{frame}

\begin{frame}[fragile]{Card--Krueger: \texttt{sample-restrictions-note}}
  {\small \artifact{.claude/skills/sample-restrictions-note/SKILL.md} (shape of the file)}
  \begin{lstlisting}[style=markdown,basicstyle=\ttfamily\footnotesize]
---
name: sample-restrictions-note
description: Use after any change to the Card-Krueger sample.
---
1. List every rule that drops stores.
2. Give the store count after each step.
3. Quote the command that printed the counts (panel_workflow.py).
4. Never edit the data file.
  \end{lstlisting}
  {\small The header lines between the dashes: its name and when to use it.}
\end{frame}

\begin{frame}{Your turn}
  \youtry{Thu 11:25--11:55}
  \vspace{0.3em}
  \begin{enumerate}
    \item Fill in \artifact{templates/skill\_template.md}.
    \item On your separate copy, save it as \artifact{.claude/skills/<name>/SKILL.md}.
    \item Run it once on the Card--Krueger folder.
    \item Fix the one step the agent misread.
  \end{enumerate}
\end{frame}

\begin{frame}{Thursday wrap and Friday recap}
  \timebox{Thu 11:55--12:00 \textperiodcentered\ Fri 09:30--09:35}
  \vspace{0.3em}
  \begin{block}{Tonight}
    Draft one Skill for a task you repeat every week (\artifact{docs/between\_sessions.md}).
  \end{block}
  \vspace{0.4em}
  \begin{block}{Friday 09:30: so far}
    Agent \textperiodcentered\ saved versions \textperiodcentered\
    rules file \textperiodcentered\ to-do list \textperiodcentered\
    Skills
  \end{block}
\end{frame}

% ----------------------------------------------------------------
\begin{frame}{Sub-agent: what it is}
  \timebox{Fri 09:35--09:55}
  \begin{block}{In plain words}
    A sub-agent is another agent you hand a piece of the work. It has its own
    conversation and its own instructions. Inside the bounds you set, it can read,
    write, check, and decide. It hands back whatever that job needs: a cleaned file,
    a figure, a review, a pull request, or a short report. A summary is one possible
    result, not the definition.

    A Skill is a repeatable set of instructions followed in the same conversation.
    A sub-agent carries out the piece of work on its own, in a fresh context.
  \end{block}
  \vspace{0.4em}
  {\textcolor{courseteal}{\textbf{Think of it as:}} a research assistant who redoes
    the Card--Krueger sample, rebuilds the parallel-trends figure, or reviews the
    cleaning, and hands the result back.}
\end{frame}

\begin{frame}{Card--Krueger: Skill vs sub-agent}
  \small
  \begin{tabular}{@{}p{0.18\textwidth}p{0.38\textwidth}p{0.38\textwidth}@{}}
    \toprule
     & \textbf{Skill} & \textbf{Sub-agent} \\
    \midrule
    What it is & Repeatable instructions, in this conversation & Another agent, with its own instructions \\
    Where & Same conversation & Its own conversation (a fresh context) \\
    Hands back & The steps, done in place & A file, a figure, a review, a pull request, or a short report \\
    Example & \artifact{replication-checker} & \artifact{data-reviewer}, or the parallel-trends figure \\
    \bottomrule
  \end{tabular}

  \vspace{0.9em}
  Inside the bounds you set, it can read, write, check, and decide.
  A summary is one possible result, not the definition.

  \vspace{0.35em}
  \artifact{data-reviewer} is bounded to a review of panel balance and NJ/PA coding: no edits.
\end{frame}

\begin{frame}{Your turn}
  \youtry{Fri 09:35--09:55}
  \vspace{0.3em}
  Open \artifact{.claude/agents/data-reviewer.md} and mark three things:
  \begin{enumerate}
    \item \textbf{The job:} what piece of work is this?
    \item \textbf{The limits:} what bounds did we set?
    \item \textbf{What it hands back:} a review. Name one other result a sub-agent can hand back.
  \end{enumerate}
\end{frame}

% ----------------------------------------------------------------
\begin{frame}{Rule of thumb: Skill or sub-agent?}
  \timebox{Fri 09:55--10:15}
  \inplainwords{Use a \textbf{Skill} when the same steps should be reused in place.\\
    Use a \textbf{sub-agent} when the piece of work needs its own focus, its own
    files, or would crowd the main conversation.}
  \likethis{the sample-restrictions checklist in this chat, versus a research
    assistant who redoes the sample, rebuilds the figure, or reviews the cleaning.}
\end{frame}

\begin{frame}{Card--Krueger: two tasks}
  \begin{columns}[T,onlytextwidth]
    \begin{column}{0.47\textwidth}
      \begin{block}{Skill}
        ``Recompute the DiD the same way each time.''
      \end{block}
    \end{column}
    \begin{column}{0.47\textwidth}
      \begin{block}{Sub-agent}
        ``Redo the sample, or rebuild the parallel-trends figure, and hand back the file.''
      \end{block}
    \end{column}
  \end{columns}
  \vspace{0.7em}
  You set the bounds. Inside them, the sub-agent can decide. You still approve the merge.
\end{frame}

% Answer key: 1 Skill, 2 sub-agent, 3 Skill, 4 sub-agent, 5 Skill, 6 sub-agent.
\begin{frame}{Your turn: sort six tasks}
  \youtry{Fri 09:55--10:15 \textperiodcentered\ then compare with a neighbour}
  \vspace{0.2em}
  \small
  \begin{tabular}{@{}rp{0.68\linewidth}l@{}}
    \toprule
    & \textbf{Card--Krueger task} & \textbf{Skill or sub-agent?} \\
    \midrule
    1 & Recompute group means after every data change & \\
    2 & Check NJ/PA coding without seeing my code & \\
    3 & Write the sample-restrictions note the same way & \\
    4 & Review literature claims while I fix code & \\
    5 & Check the folder reproduces from a fresh start & \\
    6 & Review my pull request in its own conversation & \\
    \bottomrule
  \end{tabular}
\end{frame}
